%% sections/06-conclusion.tex

\section{Conclusion}
\label{sec:conclusion}

We have presented the Marathon Playtest Rubric, an 8-dimension scoring instrument for
AI-simulated tabletop role-playing game sessions, and reported its application across
49 sessions spanning 7 complete campaigns with distinct party archetypes and dramatic
questions. The rubric's four core design commitments --- evidence-basis (every score
requires specific log citations), dimension independence (each dimension measures a
distinct quality), forward-only amendment (prior session scores are never retroactively
adjusted), and conservative thresholds (the 56/80 binding threshold is set below the
mature session mean) --- together produce a measurement instrument that is both stable
enough for longitudinal comparison and adaptive enough to learn from novel session
outcomes.

The central empirical findings are: a 12-point improvement in mean gate scores from
early corpus sessions (65/80) to mature sessions (77/80), driven by a combination of
genuine design improvement and 13 rubric amendments; cross-campaign consistency within
0.7 points across campaigns with fundamentally different party archetypes; and a
100\% binding-session PASS rate at the 56/80 threshold. The amendment rate declined
from 5 amendments in Campaign 1 to zero amendments in Campaigns 6 and 7, suggesting
convergence toward a stable quality description for the current corpus.

The rubric's amendment mechanism is itself a contribution independent of the specific
anchors it has produced. Any property-based scoring instrument applied to a creative
output corpus will encounter outcomes that exceed its predictive scope; the systematic
pipeline for logging those outcomes, clustering them by dimension, and codifying them
as new anchors provides a principled method for closing that gap without introducing
retrospective bias. This mechanism is applicable beyond TRPG research, to any domain
where AI-generated creative outputs require systematic quality evaluation at scale.

\subsection{Future Work}

Three priorities follow directly from the limitations identified in
Section~\ref{sec:discussion}.

First, inter-rater reliability must be established. A 10-session subset of the corpus
should be independently scored by two or more human TRPG practitioners following the
rubric's evidence-citation protocol, and Cohen's $\kappa$ or Krippendorff's $\alpha$
computed per dimension. This will establish whether the rubric's band descriptions are
precise enough to enable consistent scoring between raters and will identify dimensions
that require anchor clarification.

Second, multi-system validation is needed to test whether the rubric's anchors are
calibrated to patterns specific to a single AI system or generalize across AI DM
implementations. Running a subset of the Marathon adventures with a different LLM-based
DM and scoring the resulting sessions with the current rubric will reveal any
system-specific bias in the anchor language.

Third, extension to non-D\&D systems would expand the rubric's applicability
significantly. Adapting the instrument for Powered by the Apocalypse games, which use
a fundamentally different mechanical resolution structure, would require systematic
revision of the Mechanical Fairness and Table Readiness dimensions and would test
whether the rubric's core structure (8-dimensional, property-based, amendment-driven)
transfers to game systems with different design assumptions about player choice, risk,
and narrative authority.

A longer-term goal is a human DM comparison: running the same compiled adventure with
both an AI DM and a human DM, scoring both sessions with the same rubric, and analyzing
where the two sessions diverge on individual dimensions. Such a study would not
adjudicate whether AI DMs are ``better'' or ``worse'' --- the dimensions measure session
execution quality, not the subjective experience of the participants --- but it would
map the specific quality dimensions where AI-simulated and human-run sessions differ
systematically, which is a prerequisite for improving AI DM systems toward human-level
session quality.
